\documentclass[aspectratio=169,11pt]{beamer}

\usetheme{Madrid}
\setbeamertemplate{navigation symbols}{}

\usepackage{amsmath,amssymb}
\usepackage{graphicx}
\usepackage{currfile}

\title[Planck Polarization Separation]{Planck Polarization Component Separation in STL}
\author{Alexandros Tsouros}
\date{}

\graphicspath{{\currfiledir}{\currfiledir figures/}{figures/}}

\begin{document}

\begin{frame}[fragile]
  \titlepage
\end{frame}

\begin{frame}[fragile]{Problem setup}
We observe a Planck polarization patch:
\[
  d_Q,\ d_U\in\mathbb{R}^{H\times W}.
\]

\vspace{0.25em}
The goal is to recover two unknown signal maps
\[
  \tilde s=(\tilde s_Q,\tilde s_U),
\]
using:
\begin{itemize}
  \item the observed Q/U maps,
  \item paired nuisance samples $(n_Q^{(i)},n_U^{(i)})$,
  \item a fixed auxiliary dust tracer $d_I$ from I857.
\end{itemize}
\end{frame}

\begin{frame}[fragile]{What is optimized?}
The optimization variable is the two-map tensor
\[
  u=\begin{bmatrix}u_Q\\u_U\end{bmatrix}\in\mathbb{R}^{2\times H\times W}.
\]

\vspace{0.25em}
Everything else is fixed during a minibatch step:
\[
  d_Q,\ d_U,\ d_I,\quad \{(n_Q^{(i)},n_U^{(i)})\}_{i\in I}.
\]

\vspace{0.25em}
The code initializes $u$ either from the data maps or from white noise with matching standard deviations.
After optimization, the final recovered maps are denoted
\[
  \tilde s_Q,\ \tilde s_U.
\]
\end{frame}

\begin{frame}[fragile]{Core idea: match scattering statistics}
Define a feature map $\Phi(\cdot)$ from STL scattering statistics.

\vspace{0.35em}
For a minibatch of nuisance indices $I$, compare:
\[
  y(I)=\Phi\!\left(X_{\mathrm{target}}(I)\right),
  \qquad
  \tilde y(u;I)=\Phi\!\left(X_{\mathrm{run}}(u;I)\right).
\]

\vspace{0.25em}
The loss is simply
\[
  \mathcal{L}(u;I)=\|\tilde y(u;I)-y(I)\|_2^2.
\]
\end{frame}

\begin{frame}[fragile]{Choosing noise realizations}
The nuisance files are paired before optimization:
\[
  (n_Q^{(j)},n_U^{(j)})\quad\text{share the same noise/CMB-residual seed key.}
\]

\vspace{0.25em}
At one optimization step, the script samples a set of indices
\[
  I=\{i_1,\ldots,i_{N_b}\},
\]

\vspace{0.25em}
For those indices it loads
\[
  n_Q^I,\ n_U^I\in\mathbb{R}^{N_b\times H\times W}.
\]
This leading dimension is the noise-realization dimension.
\end{frame}

\begin{frame}[fragile]{The batch dimension}
Each selected realization is treated as one batch element.
The maps that do not depend on the realization are broadcast along this dimension:
\[
  d_Q,d_U,d_I,u_Q,u_U\in\mathbb{R}^{H\times W}
  \quad\longrightarrow\quad
  \mathbb{R}^{N_b\times H\times W}.
\]

\vspace{0.25em}
After broadcasting fixed maps and inserting the sampled nuisance maps,
the channels are stacked:
\[
  X_{\mathrm{target}}(I),\ X_{\mathrm{run}}(u;I)
  \in\mathbb{R}^{N_b\times C\times H\times W}.
\]

\vspace{0.25em}
Here $C=5$ in the full phase and $C=3$ in the optional warm-up phase.
\end{frame}

\begin{frame}[fragile]{Averaging over realizations}
STL is applied to the whole batch tensor, not to one realization at a time:
\[
  X(I)\in\mathbb{R}^{N_b\times C\times H\times W}
  \quad\longmapsto\quad
  \Phi(X(I)).
\]

\vspace{0.25em}
The flattened statistics are then averaged along the batch dimension.
So the loss compares the \emph{mean} STL statistics over the selected realizations:
\[
  \bar\Phi_{\mathrm{target}}(I)
  \quad\text{against}\quad
  \bar\Phi_{\mathrm{run}}(u;I).
\]

\vspace{0.25em}
This means one optimizer step asks the same running maps $u_Q,u_U$ to work across several nuisance realizations at once.
\end{frame}

\begin{frame}[fragile]{The target channels}
For each selected nuisance pair $(n_Q^{(i)},n_U^{(i)})$, the target batch element is
\[
  X_{\mathrm{target}}^{(i)}
  =
  \begin{bmatrix}
    d_Q\\
    n_Q^{(i)}\\
    d_U\\
    n_U^{(i)}\\
    d_I
\end{bmatrix}
  \in \mathbb{R}^{5\times H\times W}.
\]

\vspace{0.25em}
Across all selected indices:
\[
  X_{\mathrm{target}}(I)\in\mathbb{R}^{N_b\times5\times H\times W}.
\]
The observed maps and auxiliary map are repeated across the batch; only the nuisance maps vary with $i$.
\end{frame}

\begin{frame}[fragile]{The running channels}
For the current running maps $u=(u_Q,u_U)$, each running batch element is
\[
  X_{\mathrm{run}}^{(i)}(u)
  =
  \begin{bmatrix}
    u_Q+n_Q^{(i)}\\
    d_Q-u_Q\\
    u_U+n_U^{(i)}\\
    d_U-u_U\\
    d_I
  \end{bmatrix}.
\]

\vspace{0.25em}
Across all selected noise realizations:
\[
  X_{\mathrm{run}}(u;I)\in\mathbb{R}^{N_b\times5\times H\times W}.
\]

\vspace{0.25em}
Each polarization field is split into:
\begin{itemize}
  \item a \emph{data-like} channel: running map plus sampled nuisance,
  \item a \emph{residual} channel: what remains in the observed map.
\end{itemize}
\end{frame}

\begin{frame}[fragile]{How fields are put into one channel}
The first and third running channels are the key construction:
\[
  u_Q+n_Q^{(i)},\qquad
  u_U+n_U^{(i)}.
\]

\vspace{0.25em}
This deliberately makes one channel look like an observed polarization map:
\[
  \text{structured signal field} + \text{one nuisance realization}.
\]

\vspace{0.25em}
The optimizer changes the same $u_Q,u_U$ for all batch elements.
The only difference between batch elements is which nuisance realization is added.
\end{frame}

\begin{frame}[fragile]{The residual channels}
The second and fourth channels are not sampled:
\[
  r_Q=d_Q-u_Q,\qquad r_U=d_U-u_U.
\]

\vspace{0.25em}
They are repeated across the minibatch:
\[
  r_Q,\ r_U \in \mathbb{R}^{H\times W}
  \quad\longrightarrow\quad
  \mathbb{R}^{N_b\times H\times W}.
\]

\vspace{0.25em}
This lets STL see whether the part removed from the data behaves like the nuisance channels.
\end{frame}

\begin{frame}[fragile]{Auxiliary dust tracer}
The I857 map is a fixed auxiliary channel:
\[
  d_I = I857 - \langle I857\rangle.
\]

\vspace{0.25em}
It is included in both target and running tensors, unchanged:
\[
  X_{\mathrm{target}}:\ d_I,\qquad X_{\mathrm{run}}:\ d_I.
\]

\vspace{0.25em}
Its role is to give the scattering statistics access to morphology correlated with polarized dust.
\end{frame}

\begin{frame}[fragile]{Which channel relations are measured?}
The 5-channel cross matrix keeps only selected statistics:
\[
\begin{pmatrix}
1&0&1&0&1\\
0&1&0&0&0\\
0&0&1&0&1\\
0&0&0&1&0\\
0&0&0&0&0
\end{pmatrix}.
\]

\vspace{0.25em}
So STL measures:
\begin{itemize}
  \item auto-statistics of all channels,
  \item Q--aux and U--aux relations,
  \item no explicit Q--U or signal--residual cross terms.
\end{itemize}
\end{frame}

\begin{frame}[fragile]{Optional warm-up without explicit noise channels}
The script can begin with a simpler 3-channel phase:
\[
  X_{\mathrm{target}}^{(i)}=
  \begin{bmatrix}d_Q\\d_U\\d_I\end{bmatrix},
  \qquad
  X_{\mathrm{run}}^{(i)}(u)=
  \begin{bmatrix}u_Q+n_Q^{(i)}\\u_U+n_U^{(i)}\\d_I\end{bmatrix}.
\]

\vspace{0.25em}
This phase matches data-like Q/U channels before introducing explicit residual channels.
\end{frame}

\begin{frame}[fragile]{Normalization and operators}
Two STL operators are built:
\[
  \Phi_3 \quad\text{for 3 channels},\qquad
  \Phi_5 \quad\text{for 5 channels}.
\]

\vspace{0.25em}
Each stores a fixed normalization reference once, then reuses it during optimization:
\[
  \texttt{store\_ref}\quad\rightarrow\quad\texttt{load\_ref}.
\]
\end{frame}

\begin{frame}[fragile]{Minibatch loop}
At each optimization step:
\begin{enumerate}
  \item choose paired nuisance indices $I$,
  \item build target statistics once,
  \item build running statistics inside the closure,
  \item take one L-BFGS optimizer step.
\end{enumerate}

\vspace{0.25em}
The target statistics are computed without gradients.
The running statistics are computed inside the closure so gradients flow back to $u_Q,u_U$.
\end{frame}

\begin{frame}[fragile]{Takeaways}
\begin{itemize}
  \item The optimized running maps are $u_Q,u_U$; the final maps are denoted $\tilde s_Q,\tilde s_U$.
  \item Target channels contain observed Q/U, sampled Q/U nuisance, and fixed I857.
  \item Running channels mix fields as $u_Q+n_Q$ and $u_U+n_U$.
  \item Residual channels $d_Q-u_Q$ and $d_U-u_U$ expose what the running maps leave behind.
  \item A minibatch is an extra leading dimension over selected nuisance realizations.
  \item The cross matrix decides which channel relations STL is allowed to match.
\end{itemize}
\end{frame}

\section{Power-spectrum constraint}

\begin{frame}[fragile]{Starting point: the object shape}
In the full phase, both target and running objects have shape
\[
  X_{\mathrm{target}}(I),\quad X_{\mathrm{run}}(u;I)
  \in\mathbb{R}^{N_b\times C\times H\times W}.
\]

\vspace{0.25em}
For the run that crashes,
\[
  N_b=15,\qquad C=5,\qquad H=W=384.
\]

\vspace{0.25em}
So one real batch tensor contains
\[
  15\times 5\times384\times384
  =
  11{,}059{,}200
\]
numbers.

\vspace{0.25em}
The factor $8$ below is because one \texttt{float64} number uses $8$ bytes:
\[
  11{,}059{,}200\ \text{numbers}\times 8\ \text{bytes/number}
  \approx 88\text{ MB}
  \quad(\approx84\text{ MiB}).
\]
\end{frame}

\begin{frame}[fragile]{Two possible objectives}
The difference is what statistics are requested from this same batch tensor:
\[
  \Phi(X)
  =
  \begin{cases}
    \Phi_{\mathrm{ST}}(X), & \texttt{ST\_COMPUTE\_PS=False},\\[0.35em]
    \bigl(\Phi_{\mathrm{ST}}(X),\,P(X)\bigr), & \texttt{ST\_COMPUTE\_PS=True}.
  \end{cases}
\]

\vspace{0.25em}
Here $\Phi_{\mathrm{ST}}$ contains mean, variance, and scattering coefficients.
The optional $P$ contains cross-power-spectrum coefficients.

\vspace{0.25em}
The final loss always compares target and running statistics:
\[
  \mathcal{L}(u;I)
  =
  \|\Phi(X_{\mathrm{run}}(u;I))-\Phi(X_{\mathrm{target}}(I))\|_2^2.
\]
\end{frame}

\begin{frame}[fragile]{Case 1: without the PS constraint}
With
\[
  \texttt{ST\_COMPUTE\_PS=False},
\]
the target side computes
\[
  y_{\mathrm{target}}
  =
  \bar\Phi_{\mathrm{ST}}\!\left(X_{\mathrm{target}}(I)\right),
\]
and the running side computes
\[
  y_{\mathrm{run}}(u)
  =
  \bar\Phi_{\mathrm{ST}}\!\left(X_{\mathrm{run}}(u;I)\right).
\]

\vspace{0.25em}
The optimized loss is
\[
  \mathcal{L}_{\mathrm{ST}}(u;I)
  =
  \|y_{\mathrm{run}}(u)-y_{\mathrm{target}}\|_2^2.
\]
\end{frame}

\begin{frame}[fragile]{Without PS: what is kept for gradients?}
For the target:
\[
  y_{\mathrm{target}}
  =
  \bar\Phi_{\mathrm{ST}}(X_{\mathrm{target}}(I))
\]
is computed inside \texttt{torch.no\_grad()}, so it is just a fixed tensor.

\vspace{0.25em}
For the running side, PyTorch keeps the graph needed for:
\begin{itemize}
  \item the full running batch object $X_{\mathrm{run}}(u;I)$, not one selected nuisance draw,
  \item the wavelet/scattering intermediate tensors,
  \item the final flattened scattering vector $y_{\mathrm{run}}(u)$.
\end{itemize}

\vspace{0.25em}
Here $I=\{i_1,\ldots,i_{N_b}\}$ contains all nuisance draws in the minibatch.
So the first axis of $X_{\mathrm{run}}$ stores all $N_b$ realizations at once.

\vspace{0.25em}
No power-spectrum tensors are formed or stored for backward.
\end{frame}

\begin{frame}[fragile]{Without PS: memory scale}
The full running object is
\[
  X_{\mathrm{run}}(u;I)
  \in\mathbb{R}^{15\times5\times384\times384}.
\]

\vspace{0.25em}
Since it is real \texttt{float64}, each entry costs $8$ bytes:
\[
  15\times5\times384\times384\times8
  \approx 88\text{ MB}.
\]

\vspace{0.25em}
If an intermediate representation of the same $X$ is complex, each entry costs $16$ bytes:
\[
  15\times5\times384\times384\times16
  \approx 177\text{ MB}.
\]

\vspace{0.25em}
The scattering transform then keeps wavelet/scattering intermediates for backward.
Those are substantial, but they are tied to the scattering operator:
\[
  \text{roughly scales with }N_b,\ C,\ H,\ W,\ J,\ L.
\]

\vspace{0.25em}
Crucially, without PS there is no extra tensor with shape
\[
  N_b\times C\times C\times N_{\mathrm{bins}}\times H\times W.
\]
\end{frame}

\begin{frame}[fragile]{Case 2: with the PS constraint}
With
\[
  \texttt{ST\_COMPUTE\_PS=True},
\]
the target side computes two blocks:
\[
  y_{\mathrm{target}}
  =
  \left(
    \bar\Phi_{\mathrm{ST}}(X_{\mathrm{target}}(I)),
    \bar P(X_{\mathrm{target}}(I))
  \right).
\]

\vspace{0.25em}
The running side computes the matching two blocks:
\[
  y_{\mathrm{run}}(u)
  =
  \left(
    \bar\Phi_{\mathrm{ST}}(X_{\mathrm{run}}(u;I)),
    \bar P(X_{\mathrm{run}}(u;I))
  \right).
\]

\vspace{0.25em}
The loss is still one squared difference:
\[
  \mathcal{L}(u;I)
  =
  \|y_{\mathrm{run}}(u)-y_{\mathrm{target}}\|_2^2.
\]
\end{frame}

\begin{frame}[fragile]{What is kept in memory with PS?}
The target PS block is again fixed:
\[
  \bar P(X_{\mathrm{target}}(I))
\]
because it is computed under \texttt{torch.no\_grad()}.

\vspace{0.25em}
The running PS block is computed from the full running batch:
\[
  X_{\mathrm{run}}(u;I)
  \in\mathbb{R}^{N_b\times C\times H\times W}
  \quad\longmapsto\quad
  \bar P(X_{\mathrm{run}}(u;I)).
\]
This block depends on $u_Q,u_U$, so PyTorch must keep its graph.

\vspace{0.25em}
This does not mean one nuisance term is kept separately.
The whole minibatch object $X_{\mathrm{run}}(u;I)$ is the input to $P$.

\vspace{0.25em}
In addition to the scattering graph, PyTorch keeps the PS intermediates needed so that
\[
  \frac{\partial \mathcal{L}_{\mathrm{PS}}}{\partial u_Q},
  \quad
  \frac{\partial \mathcal{L}_{\mathrm{PS}}}{\partial u_U}
\]
can be computed during \texttt{loss.backward()}.
\end{frame}

\begin{frame}[fragile]{With PS: the first large object}
The FFT power-spectrum operator forms binned Fourier maps:
\[
  X_{\mathrm{bin}}
  \sim
  N_b \times C \times N_{\mathrm{bins}}\times H\times W.
\]

\vspace{0.25em}
This is still built from the full $X_{\mathrm{run}}(u;I)$.
The $N_b$ axis is the full minibatch of nuisance draws.

\vspace{0.25em}
For a 384 map, the default is
\[
  N_{\mathrm{bins}}=24.
\]

\vspace{0.25em}
So, with complex dtype,
\[
  15\times5\times24\times384\times384\times16
  \approx 4.25\text{ GB}
  \quad(\approx3.95\text{ GiB}).
\]

\vspace{0.25em}
This is already much larger than the original real batch tensor.
\end{frame}

\begin{frame}[fragile]{With PS: the dominant cross object}
The PS operator then forms channel-pair products of the binned $X$ object:
\[
  X_{\mathrm{cross}}
  \sim
  N_b \times C \times C \times N_{\mathrm{bins}}\times H\times W.
\]

\vspace{0.25em}
The two $C$ factors mean ``channel against channel''.
This is why the memory scales like $C^2$, even though the final PS coefficients are later reduced.

\vspace{0.25em}
For the current run:
\[
  15\times5\times5\times24\times384\times384\times16
  \approx 21.23\text{ GB}
  \quad(\approx19.78\text{ GiB}).
\]

\vspace{0.25em}
This number matches the failed allocation:
\[
  \texttt{Tried to allocate 19.78 GiB}.
\]

\vspace{0.25em}
So the OOM is not caused by the final number of PS coefficients.
It is caused by forming this intermediate cross-product object from $X$.
\end{frame}

\begin{frame}[fragile]{With PS: what else is kept?}
For non-periodic boundaries, the FFT PS path can also build intermediate objects from the same $X$:
\begin{itemize}
  \item the binned Fourier maps $X_{\mathrm{bin}}$,
  \item inverse FFTs of those binned maps,
  \item crop masks,
  \item real-space products with the original $X$,
  \item the final reduced power-spectrum tensor $P(X)$.
\end{itemize}

\vspace{0.25em}
The final $P(X)$ tensor is small compared to $X_{\mathrm{cross}}$.
The costly part is the temporary path used to get to $P(X)$.

\vspace{0.25em}
During the running-side computation, PyTorch keeps enough of this path to compute
\[
  \frac{\partial \mathcal{L}_{\mathrm{PS}}}{\partial u_Q},
  \qquad
  \frac{\partial \mathcal{L}_{\mathrm{PS}}}{\partial u_U}.
\]
\end{frame}

\begin{frame}[fragile]{Why the coefficient count can look harmless}
The printed line
\[
  \texttt{43706/361810 finite}
\]
refers to the final flattened statistics after all reductions.

\vspace{0.25em}
It does \emph{not} count the temporary tensors used to compute those statistics:
\[
  X_{\mathrm{bin}},\quad X_{\mathrm{cross}},\quad
  \text{IFFT products},\quad \text{saved backward tensors}.
\]

\vspace{0.25em}
Thus two runs can have similar final coefficient counts:
\[
  \dim y_{\mathrm{run}}
  \quad\text{similar},
\]
but very different backward memory:
\[
  \text{scattering graph only}
  \quad\text{vs.}\quad
  \text{scattering graph + PS graph}.
\]

\vspace{0.25em}
The OOM occurs when the running PS graph is retained for differentiation.
\end{frame}

\begin{frame}[fragile]{Target side versus running side}
For both settings, target statistics are computed under
\[
  \texttt{torch.no\_grad()}.
\]

\vspace{0.25em}
So the target PS computation may temporarily allocate large tensors, but it does not keep a backward graph.

\vspace{0.25em}
The running side is different:
\[
  X_{\mathrm{run}}(u;I)
  \quad\text{depends on}\quad
  u_Q,\ u_U.
\]

\vspace{0.25em}
Therefore, with PS on, the large PS intermediates are saved for \texttt{loss.backward()}.
That is why the crash happens after the coefficient count is printed.
\end{frame}

\begin{frame}[fragile]{Summary of the comparison}
\begin{columns}[T,totalwidth=\textwidth]
\begin{column}{0.48\textwidth}
\textbf{Without PS}
\begin{itemize}
  \item compute $\Phi_{\mathrm{ST}}$,
  \item keep scattering intermediates for running side,
  \item no $P(X)$ tensors,
  \item no $N_b C^2 N_{\mathrm{bins}}HW$ object.
\end{itemize}
\end{column}
\begin{column}{0.48\textwidth}
\textbf{With PS}
\begin{itemize}
  \item compute $\Phi_{\mathrm{ST}}$ and $P$,
  \item keep scattering intermediates,
  \item keep binned FFT maps,
  \item keep cross-spectrum object,
  \item about $19.78$ GiB for one dominant tensor here.
\end{itemize}
\end{column}
\end{columns}

\vspace{0.25em}
The dominant PS memory scales as
\[
  N_b\,C^2\,N_{\mathrm{bins}}\,H\,W.
\]
\end{frame}

\begin{frame}{Results}
\centering
\IfFileExists{figures/loss_curve_planck.png}{%
  \includegraphics[width=\textwidth,height=0.82\textheight,keepaspectratio]{figures/loss_curve_planck.png}%
}{%
  \vspace{1em}
  \Large Loss curve is written to \texttt{results/loss\_curve\_planck.png}
}
\end{frame}

\end{document}
